\answer{ex:17:1}{$0.2$~W contra $81$~\textmu W; o fluxo bruto precisaria de no máximo $50$~pJ/bit.}
\answer{ex:17:2}{Com $c=0.1$: $5.6$, $8.7$, $9.2$~bit/s, limite $9.3$~bit/s; com $c=0$ e $N=1000$: $65$~bit/s.}
\answer{ex:17:3}{$B=\log_2N+P\log_2P+(1-P)\log_2\frac{1-P}{N-1}=4.87$, $4.37$, $3.29$~bits por caractere, isto é, $7.3$, $6.6$, $4.9$~bit/s; para $10$~bit/s são necessários $2.3$~caracteres por segundo.}
\answer{ex:17:4}{Variância do erro $2b/(Nm^2T)+\sigma_\xi^2$ por componente; as contribuições se igualam com $N\approx90$; limite $\tfrac12\log_2(1+0.25/0.04)\approx1.4$~bit por componente por janela.}
\answer{ex:17:5}{O par converge para $\eta_bd^2+\eta_db^2<2$, isto é, aproximadamente para $\eta<1$: com $0.8$ converge, com $1.1$ há oscilações estáveis de período $2$, com $1.5$ não periódicas, com $2$ diverge; as velocidades dos aprendizes devem ser separadas.}
\answer{ex:17:6}{Limiar $\approx4.4\sigma$; $102$~kbit/s, $0.10$~mW contra $205$~mW para o fluxo bruto, $2000$ vezes menos; só $5.7\cdot10^{-5}$ das amostras saturam (mais de $3$ disparos).}
\answer{ex:17:7}{Questão aberta. Cerca de $46$~bit/s com $\text{SNR}\approx0.9$ e cerca de $330$~bit/s com ruído independente. Se o obstáculo é o ruído parecido com o sinal, a informação extraída satura com o aumento do número de canais; se é o desconhecimento do código, ela cresce com um decodificador melhor (por exemplo, um que leve em conta os tempos dos disparos).}
